\BourbakiAppendixExercises{Exercises}

\begin{enumerate}
\def\labelenumi{\arabic{enumi})}
\tightlist
\item
  Let \(\mathrm{A}\) be a \(k\) -algebra.
\end{enumerate}

\begin{enumerate}
\def\labelenumi{\alph{enumi})}
\item
  Prove that the only regular (left or right) ideal containing AA is A.
\item
  A two-sided ideal \(\mathfrak{a}\) of \(A\) is regular as a left ideal and as a right ideal if and only if the \(k\) -algebra \(A / \mathfrak{a}\) is unital.
\end{enumerate}

\begin{enumerate}
\def\labelenumi{\arabic{enumi})}
\setcounter{enumi}{1}
\tightlist
\item
  Let \(\mathbf{A}\) be a \(k\) -algebra, and let \(\mathfrak{a}\) and \(\mathfrak{b}\) be regular (left) ideals of \(\mathbf{A}\) .
\end{enumerate}

\begin{enumerate}
\def\labelenumi{\alph{enumi})}
\item
  The ideal \(a \cap b\) is regular if and only if there exist right units u modulo a and v modulo b such that \(u - v \in a + b\) .
\item
  Deduce that \(a \cap b\) is regular if a is two-sided and right regular, or if a is maximal. c) Let K be a commutative field, L be the free associative K-algebra in two variables X and Y (III, §2, No.~7, p.~448), and A be the K-subalgebra of L consisting of the elements of degree \(\geqslant 1\) . Prove that the ideals \(A(1 - X)\) and \(A(1 - Y)\) of A are regular but that their intersection is reduced to 0.
\end{enumerate}

\begin{enumerate}
\def\labelenumi{\arabic{enumi})}
\setcounter{enumi}{2}
\tightlist
\item
  Let \(\mathbf{A}\) be a \(k\) -algebra, \(\mathfrak{a}\) a regular left ideal of \(\mathbf{A}\) , and \(u\) and \(v\) right units modulo \(\mathfrak{a}\) . \(a)\) Give an example where \(u - v\) does not belong to \(\mathfrak{a}\) (cf.~III, §9, p.~644, Exercise 1). Deduce an example where the mapping \(\widetilde{\mathfrak{a}} \mapsto \widetilde{\mathfrak{a}} \cap \mathbf{A}\) (Proposition 3 of VIII, p.~437) is not injective.
\end{enumerate}

\begin{enumerate}
\def\labelenumi{\alph{enumi})}
\setcounter{enumi}{1}
\tightlist
\item
  Suppose that the ideal \(\mathfrak{a}\) is two-sided and that the \(k\) -algebra \(B = A / \mathfrak{a}\) does not contain any nonzero two-sided ideal \(\mathfrak{b}\) such that \(B\mathfrak{b} = 0\) . Prove that \(u - v\) belongs to \(\mathfrak{a}\) . This is, in particular, the case when \(A\) is commutative.
\end{enumerate}

\begin{enumerate}
\def\labelenumi{\arabic{enumi})}
\setcounter{enumi}{3}
\tightlist
\item
  Let A be a \(k\) -algebra and \(\mathfrak{a}\) a left ideal of A; we view \(\mathfrak{a}\) as a \(k\) -algebra. Prove that we have \(\Re(A) \cap \mathfrak{a} \subset \Re(\mathfrak{a})\) (apply Proposition 7 of VIII, p.~441) and that we have equality when \(\mathfrak{a}\) is two-sided (observe that every simple pseudomodule over A is simple over \(\mathfrak{a}\) or annihilated by \(\mathfrak{a}\) ). Give an example of a left ideal \(\mathfrak{a}\) such that \(\Re(A) \cap \mathfrak{a} \neq \Re(\mathfrak{a})\) (take A to be a matrix algebra).
\end{enumerate}

¶ 5) Suppose that every pseudomodule M over A is the direct sum of AM and the sub-pseudomodule of M consisting of the elements annihilated by A. Prove that A is unital. (Taking \(M = \widetilde{A}_{s}\) , prove that A admits a right unit element u. Taking \(M = A_{s}\) , prove that we have \(Ax \neq 0\) for every nonzero element x of A. Conclude that u is also a left unit element.)

\begin{enumerate}
\def\labelenumi{\arabic{enumi})}
\setcounter{enumi}{5}
\tightlist
\item
  We say that a k-algebra is primitive if it admits a pseudomodule that is simple and faithful (that is, whose annihilator is reduced to 0).
\end{enumerate}

\begin{enumerate}
\def\labelenumi{\alph{enumi})}
\item
  Show that if the algebra \(\tilde{\mathbf{A}}\) is primitive, then so is \(\mathbf{A}\) .
\item
  If the algebra A is primitive and unital, then the algebra \(\widetilde{A}\) is not primitive.
\item
  Suppose that k is a field and that the algebra A is primitive and not unital. Prove that \(\widetilde{A}\) is primitive (prove that if there exists a faithful pseudomodule over A that is not faithful as an \(\widetilde{A}\) -module, then A admits a right unit element u, and deduce a contradiction).
\end{enumerate}

\begin{enumerate}
\def\labelenumi{\arabic{enumi})}
\setcounter{enumi}{6}
\item
  We say that the \(k\) -algebra A is Artinian if every decreasing sequence of left ideals of A is stationary. The ring \(\widetilde{\mathrm{A}}\) is Artinian if and only if the algebra A and the ring \(k\) are Artinian. Prove that if \(k\) is a field, then an Artinian algebra A without radical is unital (observe that the algebra \(\widetilde{\mathrm{A}}\) is semisimple).
\item
  Suppose that \(k\) is a field and that \(A\) is a primitive algebra over \(k\) . Prove that if \(A\) admits a finite-dimensional nonzero left ideal over \(k\) , then the algebra \(A\) is finite-dimensional over \(k\) and consequently a simple ring (cf.~Exercise 7).
\item
  Suppose that k is a field and that A is a nilpotent k-algebra, that is, that there exists an integer n such that \(A^{n}=0\) . Prove that A is finite-dimensional over k (use induction on the least integer n such that \(A^{n}=0\) ).
\item
  Suppose that \(k\) is a field, that \(A\) is an infinite-dimensional \(k\) -algebra, and that every proper left ideal of \(A\) is finite-dimensional over \(k\) .
\end{enumerate}

\begin{enumerate}
\def\labelenumi{\alph{enumi})}
\item
  Prove that we have \(Aa = 0\) for every left ideal \(a \neq A\) (consider the homomorphism \(\gamma : A \to \operatorname{End}_{k}(a)\) ).
\item
  Let \(\mathfrak{r}\) be the sum of the maximal left ideals of A. Prove that \(\mathfrak{r}\) is different from A (use Exercise 9), that \(\mathfrak{r}\) is a two-sided ideal, and that \(A / \mathfrak{r}\) is a field (cf.~Exercise 7). c) Let \(u\) be a (left or right) unit modulo \(\mathfrak{r}\) . Prove that we have \(\mathfrak{ru} = \mathfrak{r}\) (observe that we have \(Au = A\) ).
\item
  Prove that rA is zero (consider the right annihilator of r in A and reason as in a)) and then that r is zero. Conclude that A is a field.
\end{enumerate}

¶ 11) Suppose that k is a field, that A is a k-algebra, and that every decreasing sequence of subalgebras of A is stationary.

\begin{enumerate}
\def\labelenumi{\alph{enumi})}
\item
  Prove that the algebra \(\mathbf{A}\) is Artinian and that its radical is finite-dimensional over \(k\) (apply Exercise 9).
\item
  Let S be a simple component of A/ \(\Re(A)\) (Exercise 7); it is isomorphic to a matrix algebra over a field D, with center Z. Prove that D is finite-dimensional over Z (if D contains an element that is transcendental over Z, construct an infinite strictly decreasing sequence of subfields of D; if D is algebraic and infinite-dimensional over Z, construct an infinite strictly increasing sequence of subfields of D, and consider their commutants in D).
\item
  Prove that if Z is infinite-dimensional over k, then we have S = D (consider the nilpotent subalgebras of B).
\item
  Prove that Z is an algebraic extension of k. Let p be the characteristic exponent of k and L the largest separable extension of k contained in Z. Prove that \(L \cap Z^{p}\) has finite degree over \(L^{p}\) (consider a p-basis of \(L \cap Z^{p}\) over \(L^{p}\) ).
\item
  Every decreasing sequence of subextensions of L is stationary. Give an example to show that this condition does not imply that L has finite degree over k (cf.~V, §12, p.~167, Exercise 3).
\item
  Conversely, let B be a k-algebra satisfying the conditions of a) such that every simple component of \(\mathrm{B}/\Re(\mathrm{B})\) satisfies the conditions of b), c), d), and e). Prove that every decreasing sequence of subalgebras of B is stationary.
\end{enumerate}

(Reduce to the case when B is a field, with center Z. Let \((\mathrm{D}_{n})\) be a decreasing sequence of subfields of B; we may assume that the field D generated by \(D_{n}\) and Z is independent of n.~Observe that we can write \(D = E \otimes_{L} Z\) , where L is a subfield of Z of finite degree over k and E is a field of finite degree over L whose center contains L. Prove that \(D_{n}\) is then the field generated by E and \(D_{n} \cap Z\) . We are thus reduced to the case when B is a commutative field that is a p-radical extension of k. If \((\mathrm{F}_{n})\) is a decreasing sequence of subextensions of B, consider the subfield \(k(\mathrm{F}_{n}^{p^{\infty}})\) , where \(F_{n}^{p^{\infty}} = \cap_{f} F_{n}^{p^{f}}\) .)

¶ 12) a) Suppose that k is a field, that the k-algebra A is Artinian, and that every increasing sequence of subalgebras of A is stationary. Prove that A has finite degree over k. (Reason as in Exercise 11. Note that the field \(k(\mathrm{X})\) contains an infinite strictly increasing sequence of subalgebras, using VII, §1, No.~5, p.~6, Proposition 6.) b) Prove that in the polynomial algebra \(k[X]\) , every increasing sequence of subalgebras is stationary. (Let \(f \in k[X]\) be a nonconstant polynomial and A be the k-subalgebra of \(k[X]\) generated by 1 and f.~Observe that A is a principal ideal domain and \(k[X]\) is a finitely generated A-module.)

\begin{enumerate}
\def\labelenumi{\arabic{enumi})}
\setcounter{enumi}{12}
\item
  Suppose that \(k\) is a field and that \(A\) admits a finite basis over \(k\) consisting of nilpotent elements. Prove that \(A\) is nilpotent (Exercise 9; reduce to the case when the algebra \(A\) is semisimple (cf.~Exercise 7), and then of the form \(\mathbf{M}_n(\mathrm{K})\) by extending the base field).
\item
  Let A and B be k-algebras; an (A,B)-pseudobimodule is a set E endowed with the structure of a left pseudomodule over A and the structure of a right pseudomodule over B such that the underlying k-module structures coincide. This corresponds to endowing E with the structure of a bimodule over the algebras \(\widetilde{A}\) and \(\widetilde{B}\) . A multiplier of E is a pair (S,D), where S:B→E is a homomorphism of right B-pseudomodules and D:A→E is a homomorphism of left A-pseudomodules, satisfying \(a\mathrm{S}(b)=\mathrm{D}(a)b\) for \(a\in A\) , \(b\in B\) . We denote by \(\mathcal{M}(\mathrm{E})\) the set of multipliers of E; it admits a natural k-module structure.
\end{enumerate}

\(a\) ) For every \(x\in \mathrm{E}\) , we denote by \(\mathrm{D}_x\) (resp. \(\mathrm{S}_x\) ) the mapping \(a\mapsto ax\) (resp. \(b\mapsto xb\) ). Show that the pair \(m_{x} = (\mathrm{D}_{x},\mathrm{S}_{x})\) is a multiplier of E.

\begin{enumerate}
\def\labelenumi{\alph{enumi})}
\setcounter{enumi}{1}
\tightlist
\item
  For every \(m = (\mathrm{S}, \mathrm{D}) \in \mathcal{M}(\mathrm{E})\) , \(a \in A\) , and \(b \in B\) , set \(am = m_{\mathrm{D}(a)}\) and \(mb = m_{\mathrm{S}(b)}\) . Prove that this defines on the k-module \(\mathcal{M}(\mathrm{E})\) the structure of a (A, B)-pseudobimodule and that the mapping \(x \mapsto m_x\) is a homomorphism of (A, B)-pseudobimodules.
\end{enumerate}

\begin{enumerate}
\def\labelenumi{\arabic{enumi})}
\setcounter{enumi}{14}
\tightlist
\item
  Let \(\mathrm{A}\) be a \(k\) -algebra. A multiplier of \(\mathrm{A}\) is a multiplier of the (A,A)-pseudobimodule \(\mathrm{A}\) . We denote the \(k\) -algebra of multipliers of \(\mathrm{A}\) by \(\mathcal{M}(\mathrm{A})\) .
\end{enumerate}

\begin{enumerate}
\def\labelenumi{\alph{enumi})}
\item
  Prove that the mapping \(\mu: a \mapsto m_{a}\) is a k-algebra homomorphism and that we have \(am = m_{a}m\) and \(ma = mm_{a}\) for \(a \in A, m \in \mathcal{M}(A)\) . The image of the homomorphism \(\mu\) is therefore a two-sided ideal of the ring \(\mathcal{M}(A)\) .
\item
  A two-sided ideal b of a k-algebra A is called faithful if its left and right annihilators are zero. Show that if A is faithful in A, then the homomorphism \(\mu\) is injective and the ideal \(\mu(\mathrm{A})\) is faithful in \(\mathcal{M}(\mathrm{A})\) .
\item
  Let B be a unital k-algebra containing A as a two-sided ideal. For \(b \in B\) , let \(\sigma_{b}\) and \(\delta_{b}\) be the mappings from A to A such that for \(a \in A\) , we have \(\sigma_{b}(a) = ba\) and \(\delta_{b}(a) = ab\) , and set \(\pi(b) = (\sigma_{b}, \delta_{b})\) . Show that \(\pi(b) \in \mathcal{M}(A)\) for every \(b \in B\) and that the resulting mapping \(\pi : B \to \mathcal{M}(A)\) is a ring homomorphism whose restriction to A is the mapping \(a \mapsto m_{a}\) . Show that \(\pi\) is injective if the ideal A is faithful in B.
\end{enumerate}

